\section{Finite \(p\)-Groups}

Let \(p\) be a prime. Groups of order \(p^n\) for \(n \geq 1\) are often called \(p\)-groups, which have certain properties that are interesting.

\begin{theorem}[\(p\)-Groups have Non-Trivial Centres]
    \label{thm:p-groups-non-trivial-centre}%
    Let \(G\) be a \(p\)-group. Then the centre of \(G\) is non-trivial.
\end{theorem}

\begin{proof}
    Let \(G \actson G\) by conjugation. Then the centre \(Z\pqty{G}\) is the set of all size-\(1\) orbits. The orbits in this case partition \(G\), so we can write
    \[
        \abs{G} = p^n = \abs{Z\pqty{G}} + \sum_{\abs{\ccl\pqty{g}} > 1} \abs{\ccl\pqty{g}}
    \]

    We know every conjugacy class has sizes that divide \(p^n\). If they are bigger than \(1\), then \(p\) divides the size. Therefore, \(p \divides \abs{G}\) and the sum of the conjugacy class sizes larger than \(1\).

    Therefore, \(p \divides Z\pqty{G}\), giving \(\abs{Z\pqty{G}} \geq p\).
\end{proof}

\begin{theorem}[Groups of Order \(p^2\) is Abelian]
    \label{thm:order-p2-abelian}%
    A group of size \(p^2\) is always abelian.
\end{theorem}

\begin{lemma}[Quotient by Centre being Cyclic implies Abelian]
    \label{lem:quotient-by-centre-cyclic-abelian}%
    If \(\flatfrac{G}{Z\pqty{G}}\) is cyclic, then \(G\) is abelian.
\end{lemma}

\begin{proof}
    That \(\flatfrac{G}{Z\pqty{G}}\) is cyclic implies that there is some coset \(x Z \pqty{G} \in \flatfrac{G}{Z\pqty{G}}\) that generates the whole quotient. In other words, every coset has the form \(x^m Z\pqty{G}\), for some \(m \geq 0\).

    Since every \(g \in G\) can be written as some \(g = x^m z\) for some \(m \geq 0, z \in Z\pqty{G}\).

    Let \(g = x^m z\) and \(h = x^n z'\), then
    \[
        g h = x^m z x^n z' = x^{m + n} z z' = x^{n + m} z' z = x^n z' x^m z = h g
    \]
    indeed. (Every element is a power of a fixed element times something that commutes with every element.)
\end{proof}

\begin{proof}[\proofname\ of \autoref{thm:order-p2-abelian}]
    Using \autoref{thm:p-groups-non-trivial-centre}, we have \(Z\pqty{G}\) is either order \(p\) or \(p^2\), and in the latter case we are done.

    But if \(\abs{Z\pqty{G}} = p\), then \(\flatfrac{G}{Z\pqty{G}}\) is cyclic, so \(G\) is abelian, contradicting with \(\abs{Z\pqty{G}} = p\).
    
    So \(\abs{Z\pqty{G}} = p^2\) and \(G\) is indeed abelian.
\end{proof}

\begin{theorem}[\(p\)-Groups have Subgoups of order Powers of \(p\)]
    \label{thm:p-group-subgoup-order-power-of-p}
    Let \(G\) be a \(p\)-group of size \(p^n\). Then \(G\) has a subgroup of size \(p^{k}\) for all \(0 \leq k \leq n\).
\end{theorem}

\begin{proof}
    We conduct induction on \(n\). For the \(n = 1\) case, this is trivially true, since \(G \isomorphic C_p\).

    If \(n > 1\), if \(k = 0\), take \(\Bqty{e}\), and so we are done.

    Otherwise, note that the centre \(Z\pqty{G}\) is non-trivial, and for \(x \in G\) where \(x \neq e\), we have \(\ord\pqty{x}\) is a power of \(p\). If we replace with a power of itself, we can assume it has order exactly \(p\). (Alternatively, \(p\) divides the order of \(Z\pqty{G}\) and apply Cauchy's Theorem.)

    Then, \(\angBrac{x}\) has order \(p\), and since \(x \in Z\pqty{G}\), we have \(\angBrac{x} \normalsub G\). This allows us to take the quotient \(\flatfrac{G}{\angBrac{x}}\). This is a group of size \(p^{n - 1}\), and by induction it contains some subgroup of size \(p^{k - 1}\), call it \(L \subgroup \flatfrac{G}{\angBrac{X}}\).

    But by the correspondence result (between the subgroups of \(G\) and subgroups of the quotient), we find some \(K \subgroup G\) containing \(\angBrac{x}\), such that \(\flatfrac{K}{\angBrac{x}} \isomorphic L\). This gives \(\abs{K} = p^k\) ad so we are done.
\end{proof}